\documentclass[10pt,a4paper]{amsart}
\usepackage{amsmath,amssymb,amsthm,graphicx}
\usepackage{verbatim}
\newenvironment{reference}{\comment}{\endcomment}
%\newenvironment{reference}{}{}
\newenvironment{slogan}{\comment}{\endcomment}
\newenvironment{history}{\comment}{\endcomment}

% For commutative diagrams we use Xy-pic
\usepackage[all]{xy}

% We use 2cell for 2-commutative diagrams.
\xyoption{2cell}
% Theorem environments.
%
\theoremstyle{plain}
\newtheorem{theorem}[subsection]{Theorem}
\newtheorem{proposition}[subsection]{Proposition}
\newtheorem{lemma}[subsection]{Lemma}

\theoremstyle{definition}
\newtheorem{definition}[subsection]{Definition}
\newtheorem{example}[subsection]{Example}
\newtheorem{exercise}[subsection]{Exercise}
\newtheorem{situation}[subsection]{Situation}

\theoremstyle{remark}
\newtheorem{remark}[subsection]{Remark}
\newtheorem{remarks}[subsection]{Remarks}

\numberwithin{equation}{subsection}

% Macros
%
\def\lim{\mathop{\mathrm{lim}}\nolimits}
\def\colim{\mathop{\mathrm{colim}}\nolimits}
\def\Spec{\mathop{\mathrm{Spec}}}
\def\Hom{\mathop{\mathrm{Hom}}\nolimits}
\def\Ext{\mathop{\mathrm{Ext}}\nolimits}
\def\SheafHom{\mathop{\mathcal{H}\!\mathit{om}}\nolimits}
\def\SheafExt{\mathop{\mathcal{E}\!\mathit{xt}}\nolimits}
\def\Sch{\mathit{Sch}}
\def\Mor{\mathop{\mathrm{Mor}}\nolimits}
\def\Ob{\mathop{\mathrm{Ob}}\nolimits}
\def\Sh{\mathop{\mathit{Sh}}\nolimits}
\def\NL{\mathop{N\!L}\nolimits}
\def\CH{\mathop{\mathrm{CH}}\nolimits}
\def\proetale{{pro\text{-}\acute{e}tale}}
\def\etale{{\acute{e}tale}}
\def\QCoh{\mathit{QCoh}}
\def\Ker{\mathop{\mathrm{Ker}}}
\def\Im{\mathop{\mathrm{Im}}}
\def\Coker{\mathop{\mathrm{Coker}}}
\def\Coim{\mathop{\mathrm{Coim}}}

% Boxtimes
%
\DeclareMathSymbol{\boxtimes}{\mathbin}{AMSa}{"02}

%
% Macros for moduli stacks/spaces
%
\def\QCohstack{\mathcal{QC}\!\mathit{oh}}
\def\Cohstack{\mathcal{C}\!\mathit{oh}}
\def\Spacesstack{\mathcal{S}\!\mathit{paces}}
\def\Quotfunctor{\mathrm{Quot}}
\def\Hilbfunctor{\mathrm{Hilb}}
\def\Curvesstack{\mathcal{C}\!\mathit{urves}}
\def\Polarizedstack{\mathcal{P}\!\mathit{olarized}}
\def\Complexesstack{\mathcal{C}\!\mathit{omplexes}}
% \Pic is the operator that assigns to X its picard group, usage \Pic(X)
% \Picardstack_{X/B} denotes the Picard stack of X over B
% \Picardfunctor_{X/B} denotes the Picard functor of X over B
\def\Pic{\mathop{\mathrm{Pic}}\nolimits}
\def\Picardstack{\mathcal{P}\!\mathit{ic}}
\def\Picardfunctor{\mathrm{Pic}}
\def\Deformationcategory{\mathcal{D}\!\mathit{ef}}
\usepackage{cleveref}
\crefname{originallemma}{original lemma}{original lemmas}
\Crefname{originallemma}{Original lemma}{Original lemmas}
\makeatletter
\newlabel{etale-lemma-characterize-closed-immersion}{{04XV}{source}{}{sourcepointer-1}{}}
\newlabel{etale-lemma-characterize-closed-immersion@cref}{{[originallemma][0][]04XV}{[1][0][]source}{}{}{}}
\makeatother
\crefname{originallemma}{original lemma}{original lemmas}
\Crefname{originallemma}{Original lemma}{Original lemmas}
\makeatletter
\newlabel{lemma-properties-diagonal}{{03HK}{source}{}{sourcepointer-2}{}}
\newlabel{lemma-properties-diagonal@cref}{{[originallemma][0][]03HK}{[1][0][]source}{}{}{}}
\makeatother
\crefname{originallemma}{original lemma}{original lemmas}
\Crefname{originallemma}{Original lemma}{Original lemmas}
\makeatletter
\newlabel{lemma-universally-closed-quasi-compact}{{04XW}{source}{}{sourcepointer-3}{}}
\newlabel{lemma-universally-closed-quasi-compact@cref}{{[originallemma][0][]04XW}{[1][0][]source}{}{}{}}
\makeatother
\crefname{originallemma}{original lemma}{original lemmas}
\Crefname{originallemma}{Original lemma}{Original lemmas}
\makeatletter
\newlabel{spaces-lemma-representable-transformations-property-implication}{{02YO}{source}{}{sourcepointer-4}{}}
\newlabel{spaces-lemma-representable-transformations-property-implication@cref}{{[originallemma][0][]02YO}{[1][0][]source}{}{}{}}
\makeatother
\crefname{originaldefinition}{original definition}{original definitions}
\Crefname{originaldefinition}{Original definition}{Original definitions}
\makeatletter
\newlabel{spaces-properties-definition-separated}{{03BS}{source}{}{sourcepointer-5}{}}
\newlabel{spaces-properties-definition-separated@cref}{{[originaldefinition][0][]03BS}{[1][0][]source}{}{}{}}
\makeatother
\title{Quasi-separatedness and separatedness descend along a surjective universally closed morphism of algebraic spaces}
\author{The Stacks Project Authors}
\date{stacks-project:spaces-morphisms.tex:local-snapshot-20261012; GNU FDL 1.2 or later}
\begin{document}
\maketitle
\setcounter{subsection}{1}
\noindent\textit{Editorial note.} All delivered mathematical text is the author's own, copied from the pinned Stacks Project chapter file spaces-morphisms.tex; no statement or proof step was rewritten, reordered or removed.\par\medskip
\noindent\textit{Editorial note.} The proof uses two lemmas of the same chapter that are counted separately in this lane; they are consumed as inputs inside the counted proof span and are declared as navigation pointers for the assembly.\par\medskip
\noindent\textit{Editorial note.} The five references of the counted spans are recorded as non-load-bearing navigation pointers; the argument that consumes them stays inside the counted proof span.\par\medskip
\noindent\textit{Editorial note.} Preamble: the three native excerpt spans of sources/stacks-spaces-morphisms/preamble.tex (lines 8-13, 14-18 and 154-217). The chapter file itself carries no preamble text.\par\medskip
\noindent\textit{Allowed standard prerequisites (editorial).}\begin{itemize}
\item A universally closed morphism of algebraic spaces is quasi-compact. Lemma lemma-universally-closed-quasi-compact, used at proof lines 1806-1807. Makes the right vertical arrow of the diagonal square quasi-compact.
\item The diagonal is representable, locally of finite type and a monomorphism. Lemma lemma-properties-diagonal, used at proof lines 1823-1825. Supplies the three properties needed to turn universal closedness into a closed immersion.
\item A universally closed representable morphism is a closed immersion. Etale Morphisms, Lemma etale-lemma-characterize-closed-immersion, with Spaces, Lemma spaces-lemma-representable-transformations-property-implication, used at proof lines 1826-1831. Delivers the closed immersion conclusion for Delta\_\{Y/B\}.
\item Separatedness of an algebraic space over S is separatedness relative to Spec(Z). Properties of Spaces, Definition spaces-properties-definition-separated, used at proof lines 1791-1794. Reduces parts (1) and (2) to parts (3) and (4).
\item Universally closed and surjective morphisms interact with quasi-compactness of base change. used at proof lines 1802-1815. Is the mechanism of the descent of quasi-separatedness.
\end{itemize}\medskip
\ifdefined\hypertarget\hypertarget{sourcepointer-1}{}\fi
\noindent\textit{Source reference.} Original lemma 04XV: focused-seeds/AGP/sources/stacks-snapshot-20261011/tags.txt, source lines 6408--6408. Navigation pointer into the Stacks Project; the standard fact it supplies is applied inside the counted proof span and the argument resting on it stays there.\par\medskip
\ifdefined\hypertarget\hypertarget{sourcepointer-2}{}\fi
\noindent\textit{Source reference.} Original lemma 03HK: focused-seeds/AGP/sources/stacks-snapshot-20261011/tags.txt, source lines 4570--4570. Navigation pointer into the Stacks Project; the standard fact it supplies is applied inside the counted proof span and the argument resting on it stays there.\par\medskip
\ifdefined\hypertarget\hypertarget{sourcepointer-3}{}\fi
\noindent\textit{Source reference.} Original lemma 04XW: focused-seeds/AGP/sources/stacks-snapshot-20261011/tags.txt, source lines 6409--6409. Navigation pointer into the Stacks Project; the standard fact it supplies is applied inside the counted proof span and the argument resting on it stays there.\par\medskip
\ifdefined\hypertarget\hypertarget{sourcepointer-4}{}\fi
\noindent\textit{Source reference.} Original lemma 02YO: focused-seeds/AGP/sources/stacks-snapshot-20261011/tags.txt, source lines 3883--3883. Navigation pointer into the Stacks Project; the standard fact it supplies is applied inside the counted proof span and the argument resting on it stays there.\par\medskip
\ifdefined\hypertarget\hypertarget{sourcepointer-5}{}\fi
\noindent\textit{Source reference.} Original definition 03BS: focused-seeds/AGP/sources/stacks-snapshot-20261011/tags.txt, source lines 4359--4359. Navigation pointer into the Stacks Project; the standard fact it supplies is applied inside the counted proof span and the argument resting on it stays there.\par\medskip
\begin{lemma}
\label{lemma-image-universally-closed-separated}
Let $S$ be a scheme. Let $B$ be an algebraic space over $S$.
Let $f : X \to Y$ be a surjective universally closed
morphism of algebraic spaces over $B$.
\begin{enumerate}
\item If $X$ is quasi-separated, then $Y$ is quasi-separated.
\item If $X$ is separated, then $Y$ is separated.
\item If $X$ is quasi-separated over $B$, then $Y$ is quasi-separated over $B$.
\item If $X$ is separated over $B$, then $Y$ is separated over $B$.
\end{enumerate}
\end{lemma}

\begin{proof}
Parts (1) and (2) are a consequence of (3) and (4) for
$S = B = \Spec(\mathbf{Z})$ (see
Properties of Spaces,
Definition \ref{spaces-properties-definition-separated}).
Consider the commutative diagram
$$
\xymatrix{
X \ar[d] \ar[rr]_{\Delta_{X/B}} & & X \times_B X \ar[d] \\
Y \ar[rr]^{\Delta_{Y/B}} & & Y \times_B Y
}
$$
The left vertical arrow is surjective (i.e., universally surjective).
The right vertical arrow is universally closed as a composition
of the universally closed morphisms
$X \times_B X \to X \times_B Y \to Y \times_B Y$. Hence it is also
quasi-compact, see
Lemma \ref{lemma-universally-closed-quasi-compact}.

\medskip\noindent
Assume $X$ is quasi-separated over $B$, i.e.,  $\Delta_{X/B}$ is
quasi-compact. Then if $Z$ is quasi-compact and $Z \to Y \times_B Y$ is
a morphism, then $Z \times_{Y \times_B Y} X \to Z \times_{Y \times_B Y} Y$
is surjective and $Z \times_{Y \times_B Y} X$ is quasi-compact by our remarks
above. We conclude that $\Delta_{Y/B}$ is quasi-compact, i.e., $Y$
is quasi-separated over $B$.

\medskip\noindent
Assume $X$ is separated over $B$, i.e., $\Delta_{X/B}$ is a closed
immersion. Then if $Z$ is affine, and $Z \to Y \times_B Y$ is
a morphism, then $Z \times_{Y \times_B Y} X \to Z \times_{Y \times_B Y} Y$
is surjective and $Z \times_{Y \times_B Y} X \to Z$ is universally closed
by our remarks above. We conclude that $\Delta_{Y/B}$ is universally closed.
It follows that $\Delta_{Y/B}$ is representable, locally of finite type, a
monomorphism (see
Lemma \ref{lemma-properties-diagonal})
and universally closed, hence a closed immersion, see
\'Etale Morphisms,
Lemma \ref{etale-lemma-characterize-closed-immersion}
(and also the abstract principle
Spaces, Lemma
\ref{spaces-lemma-representable-transformations-property-implication}).
Thus $Y$ is separated over $B$.
\end{proof}

\end{document}
